\section{Convergent inversion at a flat critical value}\label{sec:flat}

The key analytic problem after complete cancellation is
\begin{equation}\label{eq:flat-model}
 F(t)=C\,R(\ee^{-S/t}),\qquad
 R(Q)=1+c_\nu Q^\nu H(Q),\qquad H(0)=1,
\end{equation}
where \(C\ne0\), \(S>0\), \(c_\nu\ne0\), and \(\nu\ge1\).
The function \(R\) is holomorphic near zero.
The output approaches \(C\) faster than every power of \(t\)
on a strict right subsector.
The finite multiplicity is a property of the \emph{nome map}
\(Q\mapsto R(Q)\), not a finite Taylor multiplicity of \(F\)
at the boundary point \(t=0\).

\subsection{All coefficients and all local sheets}

Put
\[
 Y=\frac{y/C-1}{c_\nu}.
\]
Choose the analytic root \(H(Q)^{1/\nu}\) with value \(1\) at
zero.  Then \(\chi(Q)=QH(Q)^{1/\nu}\) satisfies
\(\chi(0)=0,\chi'(0)=1\), and hence has an analytic inverse
\(Q=\Psi(u)\).

\begin{theorem}[Complete small-nome flat inverse]\label{thm:flat-inverse}
There is a disk on which \(\Psi\) is holomorphic,
\(\Psi(0)=0,\Psi'(0)=1\).  Its coefficients and those of
\(\ell(u)=\log(\Psi(u)/u)\) are
\begin{align}
 \Psi(u)&=\sum_{n\ge1}\frac{u^n}{n}
             [z^{n-1}]H(z)^{-n/\nu},\label{eq:psi-lagrange}\\
 \ell(u)&=\sum_{n\ge1}\frac{u^n}{n}
             [z^n]H(z)^{-n/\nu}.\label{eq:ell-lagrange}
\end{align}
For \(Y\ne0\) sufficiently small, fix a lift \(\Log Y\).
Every inverse with sufficiently small \(|\ee^{-S/t}|\) is exactly
\begin{align}
 u_j&=\exp\left(\frac{\Log Y+2\pi\ii j}{\nu}\right),\notag\\
 t_j(Y)&=
 -\frac{S}{(\Log Y+2\pi\ii j)/\nu+\ell(u_j)},
 \qquad j\in\Z.\label{eq:all-flat-sheets}
\end{align}
These branches lie in \(\Re t>0\) when \(|Y|\) is sufficiently
small.  A positive loop about \(Y=0\) sends \(j\) to \(j+1\), and
\begin{equation}\label{eq:deck}
 t_{j+\nu}=\frac{t_j}{1-2\pi\ii t_j/S}.
\end{equation}
\end{theorem}
\begin{proof}
The analytic inverse-function theorem gives \(\Psi\).
Its equation is
\(\Psi=uH(\Psi)^{-1/\nu}\).
Lagrange--B\"urmann proves \cref{eq:psi-lagrange}.
Applying the same formula to
\(\ell(u)=-(1/\nu)\log H(\Psi(u))\), or differentiating
\(H(z)^{-n/\nu}\) inside the coefficient formula, gives
\cref{eq:ell-lagrange}.

The equation \(R(Q)=y/C\) is equivalent to
\(\chi(Q)^\nu=Y\).
It has exactly the \(\nu\) roots
\(\Psi(u_j)\), with \(j\) modulo \(\nu\), in a sufficiently
small \(Q\)-disk.  Solving
\(\ee^{-S/t}=\Psi(u_j)\) requires all logarithm lifts.
Since
\[
 \log\Psi(u_j)=
 \frac{\Log Y+2\pi\ii j}{\nu}+\ell(u_j)
\]
is already a coherent lift, allowing \(j\) to range over all
integers incorporates both the finite roots and all logarithmic
lifts exactly once.  Increasing \(j\) by \(\nu\) leaves \(u_j\)
unchanged and adds \(2\pi\ii\) to the denominator.  This proves
\cref{eq:deck} and the monodromy statement.

The real part of the denominator is
\(\log|\Psi(u_j)|<0\) for small \(Y\); hence
\(-S\) divided by it has positive real part.
Completeness is asserted in the small-nome region, where all
roots of the analytic equation have just been counted.
\end{proof}

The distinction in the last sentence matters.  The theorem does
not classify every possible tangential approach to \(t=0\) for
which \(|Q|\) remains near \(1\).
Its natural cusp region is
\(\Re(S/t)>T\), with \(T\) sufficiently large.

\subsection{The resulting transseries and its support}

On one logarithm lift, put \(L=-\log u\).
Then
\begin{equation}\label{eq:log-puiseux-transseries}
 t=\frac{S}{L-\ell(u)}
   =\frac{S}{L}
     +S\sum_{n\ge1}u^n
       \sum_{a=1}^n
        \frac{[u^n]\ell(u)^a}{L^{a+1}},
\end{equation}
for sufficiently small \(u\).  A sufficient condition for the
regrouped series is
\(\sum_{n\ge1}|[u^n]\ell(u)|\,|u|^n<|L|\).
This is a convergent logarithmic--Puiseux representation.
For each \(u^n\), the coefficient is a finite polynomial in
\(L^{-1}\), with no hidden infinite inner sum.
Since \(u=Y^{1/\nu}\) on a chosen sheet, the power scale in the
output variable is rationally ramified.

Let \(\mathcal A\) be the additive monoid generated by the
positive integers in the support of \(H-1\).
\Cref{eq:psi-lagrange,eq:ell-lagrange} show
\[
 \operatorname{supp}\ell\subseteq\mathcal A.
\]
Indeed every coefficient of \(H^{-n/\nu}\) is a finite sum of
products of coefficients of \(H-1\).  Consequently no exponent
outside \(\mathcal A\) can occur in \(\ell\).
This also shows which actions must be re-examined after
cancellation: they are the \emph{gaps above the first surviving
action}, encoded by \(H\), not a list of the ungrouped input
factors.

\subsection{An explicit uniform error bound}

\begin{theorem}[Certified logarithmic--Puiseux truncation]
\label{thm:flat-tail}
Suppose \(H\) is holomorphic on a neighborhood of \(|Q|\le r\)
and
\[
 |H(Q)-1|\le\frac12\qquad (|Q|\le r).
\]
Put
\[
 R_0=r\,2^{-1/\nu},\qquad M=\frac{\log2}{\nu}.
\]
Then \(\Psi\) and \(\ell\) are holomorphic for \(|u|<R_0\),
and \(|\ell(u)|\le M\).
If \(\ell_N\) is the Taylor polynomial through degree \(N\),
\(\rho=|u|/R_0<1\), then
\begin{equation}\label{eq:ell-tail}
 |\ell(u)-\ell_N(u)|
 \le \frac{M\rho^{N+1}}{1-\rho}.
\end{equation}
For any coherent logarithm lift \(L=-\log u\), assume
\(\rho\le1/2\) and \(|L|\ge2M\).
The approximation \(t_N=S/(L-\ell_N(u))\) satisfies
\begin{equation}\label{eq:flat-certified}
 |t-t_N|
 \le
 \frac{4SM\,\rho^{N+1}}{(1-\rho)|L|^2}.
\end{equation}
\end{theorem}
\begin{proof}
The disk \(|H-1|\le1/2\) excludes zero and defines the chosen
root and logarithm.  On \(|Q|=r\),
\[
 |QH(Q)^{1/\nu}|\ge r\,2^{-1/\nu}=R_0.
\]
Rouch\'e's theorem counts exactly one simple root of
\(QH(Q)^{1/\nu}=u\) in \(|Q|<r\) for \(|u|<R_0\).
This extends \(\Psi\) analytically over the entire \(u\)-disk.
Since
\[
 \ell(u)=-\frac1\nu\log H(\Psi(u)),
\]
the power series for \(\log(1+z)\) gives
\(|\ell|\le(\log2)/\nu=M\).
Cauchy's estimate proves \cref{eq:ell-tail}.

There is no constant term in \(\ell\).  Therefore, when
\(\rho\le1/2\),
\[
 |\ell_N(u)|\le M\sum_{n=1}^N\rho^n
                 \le\frac{M\rho}{1-\rho}\le M.
\]
Both \(L-\ell\) and \(L-\ell_N\) have modulus at least
\(|L|/2\).  Subtracting their reciprocals proves
\cref{eq:flat-certified}.
\end{proof}

For all the lifts in \cref{eq:all-flat-sheets},
\[
 \Re(-\log u_j)=\frac1\nu\log(1/|Y|).
\]
Thus the absolute error estimate can be made uniform over all
integer sheet indices in the small-nome region.
This does not say that all these points remain in one fixed
angular subsector of \(t\).  As \(|j|\) grows at fixed \(Y\),
the points approach \(0\) tangentially.
If a fixed angular subsector is required, it imposes a bound of
order \(\log(1/|Y|)\) on the admissible magnitude of \(j\).

\subsection{Flatness, criticality, and nonconstancy}

On a closed right subsector,
\(|Q|\le\ee^{-S\sin\delta/|t|}\), so
\(F(t)-C=o(|t|^N)\) for every \(N\).
A nonconstant function of this form cannot extend
holomorphically to a neighborhood of \(t=0\):
such an extension would have every Taylor coefficient of
\(F-C\) equal to zero, hence would be constant.
This explains why ordinary analytic ramification at \(t=0\)
does not describe the problem.

A nonzero finite eta exponent vector with
\(A_{h,k}=B=0\) cannot produce a constant function.
If it did, analytic continuation in the upper half-plane would
make the original eta quotient constant everywhere.
As \(q\to0\) at infinity it is
\[
 q^{D/24}\prod_dP(q^d)^{r_d}.
\]
For \(D\ne0\) this is not constant.
For \(D=0\), let \(d_*\) be the smallest index with \(r_{d_*}\ne0\).
The coefficient at \(q^{d_*}\) is \(-r_{d_*}\ne0\).
This contradiction proves that a first nonzero dual coefficient
exists in the balanced eta case.
The analogous argument applies to a nonzero raw product satisfying
all three balances.

The derivative of the inverse near the flat value is highly
sensitive.  From its leading term,
\[
 t\sim\frac{\nu S}{\log(1/Y)},\qquad
 \frac{\mathrm dt}{\mathrm dY}
 \sim\frac{\nu S}{Y\log^2(1/Y)}
\]
on a fixed sheet.
A tiny absolute error in the measured flat defect may therefore
destroy useful relative information.  Keeping the exact modular
normalization and using logarithmic variables are part of the
mathematics of stable inversion.
